\documentclass[11pt]{article}
\usepackage[margin=1in]{geometry}
\usepackage{amsmath,amssymb,amsthm,mathtools}
\usepackage{booktabs,longtable,array}
\usepackage{seqsplit}
\usepackage{microtype}
\usepackage{enumitem}
\usepackage{xcolor}
\usepackage[colorlinks=true,linkcolor=blue!55!black,citecolor=blue!55!black,urlcolor=blue!55!black]{hyperref}
\usepackage[T1]{fontenc}
\usepackage{lmodern}

\newtheorem{theorem}{Theorem}[section]
\newtheorem{lemma}[theorem]{Lemma}
\newtheorem{proposition}[theorem]{Proposition}
\newtheorem{corollary}[theorem]{Corollary}
\newtheorem{definition}[theorem]{Definition}
\newtheorem{remark}[theorem]{Remark}
\newtheorem{conjecture}[theorem]{Conjecture}
\newtheorem{firewall}[theorem]{Firewall}

\newcommand{\R}{\mathbb R}
\newcommand{\C}{\mathbb C}
\newcommand{\N}{\mathbb N}
\newcommand{\1}{\mathbf 1}
\newcommand{\eps}{\varepsilon}
\newcommand{\OO}{\mathcal O}
\newcommand{\PP}{\mathcal P}
\newcommand{\TT}{\mathcal T}
\newcommand{\RR}{\mathcal R}
\newcommand{\SSS}{\mathcal S}
\newcommand{\dd}{\,d}
\newcommand{\zetaR}{\zeta}
\newcommand{\GPHT}{\mathrm{GPHT}_{2,3}}
\newcommand{\code}[1]{\texttt{#1}}
\newcommand{\sha}[1]{\texttt{\seqsplit{#1}}}

\title{Parity-Covariant Gluing and Fixed-Depth Obstructions for Native M\"obius Rows\\
\large A Hostile Reconstruction of a Two-Row Mellin--Landau Route}
\author{Agentic Polymath Project}
\date{17 August 2026}

\hypersetup{pdftitle={Parity-Covariant Gluing and Fixed-Depth Obstructions for Native Mobius Rows},pdfauthor={Agentic Polymath Project},pdfsubject={Hostile reconstruction of a projective paired-source route to the Riemann Hypothesis}}

\begin{document}
\maketitle

\begin{abstract}
A projective paired-source proposal attempted to prove coefficientwise
nonnegativity of a canonical M\"obius row by stopping a rough-prime recursion at
one-prime leaves, applying one canonical Target--Lorenz transport to every
leaf, and then invoking a two-row Mellin--Landau argument.  We give a standalone
hostile reconstruction of the complete chain.

The coefficient magnitudes and every component-row activation indicator are
transported exactly through arbitrary rough-prime histories.  The missing
interface is a parity cocycle: a history of length $m$ multiplies the signed
terminal marginal by $(-1)^m$.  The published terminal realization is
one-sided and is not invariant under this cocycle.  We prove a necessary and
sufficient parity-covariant gluing theorem and exhibit the explicit odd-history
leaf
\[
 X=67\cdot71\cdot13,\qquad h=(67),\qquad(p,y)=(71,13).
\]
Exact dyadic rational interval arithmetic over all $239$ active $P_{61}$
divisors proves
\[
 E_T(71,13)-O_T(71,13)>17,
\]
whereas the odd history requires the reverse Target--Lorenz orientation.

We next show that the obvious parity repair also fails.  Iterating any fixed
depth $L$ gives an exact disjoint source decomposition, but its current block
has leading sign $(-1)^{L-1}(\log\log X)^{L-1}$.  Hence an even depth restores
canonical recursive parity only by making the current block eventually
negative.  A $256$-bit MPFR certificate already gives
\[
 B_{2,200000}(2)<-11,\qquad B_{2,200000}(3)<-2.
\]
Thus no fixed-depth local reset can close the route.

As a clean replacement we prove a global parity-aware common-source gluing
theorem.  At each finite endpoint it either produces one source-faithful
positive row after all history parities are incorporated, or emits a separating
dual certificate.  Its two-row specialization $\GPHT$ is the exact remaining
global producer.  Finally, we reprove the fixed-row reciprocal-zeta Mellin
transform, a self-contained Landau lemma, and the fact that rows two and three
have no common open-strip cancellation.  Therefore $\GPHT$ implies the Riemann
Hypothesis.  The producer is not proved here.  The unconditional content of the
paper is the corrected gluing theorem, the explicit counterexample, the
fixed-depth no-go theorem, the global primal-dual architecture, and the exact
Mellin--Landau reduction.  The Riemann Hypothesis remains unproved.
\end{abstract}

\begin{center}
\fbox{\parbox{0.92\textwidth}{\textbf{Scientific status.}
This manuscript is an unconditional correction and reduction, not an
announcement of a proof of the Riemann Hypothesis.  It withdraws the parity-blind
composition used in the predecessor full-row proposal and isolates a new
parity-aware producer.}}
\end{center}

\tableofcontents

\section{Introduction}

The route studied here starts from a family of positive component rows
$Q_X(j)$ and forms the signed M\"obius row
\begin{equation}
 c_X(j)=\sum_{k\le X/j}\frac{\mu(k)}{\sqrt{k}}Q_{X/k}(j),
 \qquad j\ge2.
 \label{eq:native-row}
\end{equation}
A proposed proof of $c_X(j)\ge0$ for all $X$ and $j$ used three ideas:

\begin{enumerate}[label=(\roman*)]
\item encode the sign $\mu(k)$ as the difference of two positive parity channels;
\item expand squarefree $k$ by ordered rough-prime histories and stop when the
remaining quotient lies below $67$;
\item realize every stopped one-prime leaf by the same positive Target--Lorenz
map, then sum the resulting positive rows.
\end{enumerate}

The downstream analytic implication is unusually short.  If rows two and three
are nonnegative, their Mellin transforms contain $1/\zeta(s+1/2)$, while their
finite Mellin numerators cannot vanish simultaneously at a point of the open
critical strip.  A nonnegative-transform theorem of Landau then excludes every
zero with real part greater than $1/2$.

The burden is therefore entirely upstream: does the positive source tree really
produce the signed row in \eqref{eq:native-row} without changing signs, copying
source atoms, or dropping activations?  This paper reconstructs that interface
at atom level.

The answer separates cleanly into two parts.  Magnitudes and activations are
correct.  If a rough prime $p$ is moved from the source index into the path,
then
\[
 p^{-1/2}(k/p)^{-1/2}=k^{-1/2},
 \qquad
 \frac{X/p}{k/p}=\frac{X}{k}.
\]
The sign is not invariant: the parity channels are exchanged once, and the
signed observation changes sign.  After $m$ rough primes, the terminal datum is
$(-1)^m$ times its canonical local orientation.  The predecessor proof retained
this parity in its labels but omitted it from the terminal realization.

Our main unconditional contributions are:

\begin{itemize}
\item a necessary-and-sufficient parity-covariant projective gluing theorem;
\item an exact proof that rough placement preserves every row activation and
coefficient magnitude while contributing the cocycle $(-1)^{|h|}$;
\item a proof-grade odd-history obstruction with target gap greater than $17$;
\item an exact depth-$L$ source expansion and a no-go theorem for every fixed
local parity reset;
\item a global parity-aware common-source primal-dual theorem;
\item a self-contained two-row Mellin--Landau consumer.
\end{itemize}

The corrected architecture does not terminalize one-prime leaves independently
and does not attempt another fixed-depth reset.  It expands the finite source at
a fixed endpoint, incorporates every cumulative history parity, and applies one
global common-source transport.  The resulting producer is $\GPHT$.  Separate
leafwise certificates cannot prove it, because they would reuse source capacity
and erase the parity cocycle.

\section{Canonical component rows and native M\"obius rows}

\subsection{The component row}

For $j\ge2$, set
\[
 A_j=\frac{j+1}{j-1},\qquad
 B_j=\frac{(j+1)(j-2)}{j(j-1)},\qquad
 C_j=\frac{2}{j(j-1)}.
\]
The canonical component row is
\begin{align}
 Q_X(j)={}&
 \frac{A_j}{\sqrt j}\log\frac Xj\,\1_{X\ge j}
 -\frac{B_j}{\sqrt{j+1}}\log\frac X{j+1}\,\1_{X\ge j+1}
 \notag\\
 &+C_j\sum_{m\ge j+2}
 \frac1{\sqrt m}\log\frac Xm\,\1_{X\ge m}.
 \label{eq:Q}
\end{align}
Every sum in this paper is causally extended by zero.  In particular,
$Q_X(j)=0$ when $X<j$.  Between consecutive integer knots, $Q_X(j)$ is affine
in $\log X$.

The full native row is \eqref{eq:native-row}.  No sign assertion is built into
its definition.  The exact ordinary-capacity and entropy identities that
motivated the row are useful background, but the present paper uses only the
component formula \eqref{eq:Q} and its Mellin transform.

\subsection{Squarefree source notation}

Let
\[
 P_{61}=\prod_{p\le61}p.
\]
Every squarefree source index has a unique factorization
\begin{equation}
 k=d\,p_1\cdots p_t,
 \qquad d\mid P_{61},
 \qquad 67\le p_1<\cdots<p_t.
 \label{eq:factorization}
\end{equation}
We call $h(k)=(p_1,\ldots,p_t)$ its ordered rough history.  Since $k$ is
squarefree,
\begin{equation}
 \mu(k)=\mu(d)(-1)^t.
 \label{eq:mu-parity}
\end{equation}
The factor $(-1)^t$ is the cocycle omitted by the parity-blind gluing step.

\section{Positive paired sources and the parity cocycle}

\subsection{Paired source category}

A paired source is a pair $P=(E,O)$ of positive measures.  Let
\[
 S(E,O)=(O,E)
\]
be the parity swap.  Given a positive row map $R$, define signed observation by
\begin{equation}
 \OO(P)=R(E)-R(O).
 \label{eq:observation}
\end{equation}
The fundamental identity is immediate.

\begin{lemma}[Parity cocycle]
For every integer $m\ge0$,
\begin{equation}
 \OO(S^mP)=(-1)^m\OO(P).
 \label{eq:parity-cocycle}
\end{equation}
\end{lemma}

\begin{proof}
The claim for $m=1$ follows by exchanging the two terms in
\eqref{eq:observation}; induction gives the result.
\end{proof}

\subsection{Activation and coefficient covariance}

The concern that arbitrary rough histories might move a row across an activation
boundary is natural.  It is not the failing interface.

\begin{theorem}[Exact rough-placement functor]
\label{thm:activation}
Suppose $p\mid k$ is one rough factor.  Moving $p$ from the source index into
the path sends $(X,k)$ to $(X/p,k/p)$ and multiplies the child coefficient by
$p^{-1/2}$.  Then
\begin{equation}
 \frac{X/p}{k/p}=\frac Xk,
 \qquad
 p^{-1/2}(k/p)^{-1/2}=k^{-1/2}.
 \label{eq:activation-functor}
\end{equation}
Consequently every indicator $\1_{X/k\ge j}$ and every value
$k^{-1/2}Q_{X/k}(j)$ is transported without change in magnitude.  After a
history $h$ the signed value is
\begin{equation}
 (-1)^{|h|}k^{-1/2}Q_{X/k}(j).
 \label{eq:history-signed-row}
\end{equation}
\end{theorem}

\begin{proof}
Both identities in \eqref{eq:activation-functor} are algebraic.  The row
activation depends only on the ratio $X/k$.  The sign statement follows from
\eqref{eq:parity-cocycle}.
\end{proof}

\begin{remark}
The theorem covers one-sided activation boundaries exactly.  If $X/k=j$, then
both parent and placed child lie on the same boundary.  No continuity or
floating-point argument is involved.
\end{remark}

\section{The corrected projective gluing theorem}

We now state the abstract interface in the form needed for a rigorous source
composition.

\begin{definition}[Labelled projective source tree]
A labelled projective source tree consists of:
\begin{enumerate}[label=(\alph*)]
\item a finite or finitely terminating rooted tree;
\item at each node $v$, a positive paired source $P_v$ and an ordered rough
history $h(v)$;
\item a one-use direct-sum source identity expressing each nonterminal $P_v$ as
current source plus positive multiples of child sources;
\item at every terminal node, a canonical local signed datum $K_v$.
\end{enumerate}
The root-oriented terminal datum is
\[
 \eps(v)K_v,
 \qquad \eps(v)=(-1)^{|h(v)|}.
\]
\end{definition}

\begin{theorem}[Parity-covariant projective gluing]
\label{thm:parity-gluing}
Assume every source atom has one owner and every recursive coefficient is
nonnegative.  A family of positive terminal rows preserves the exact root
marginal if each terminal row realizes
\begin{equation}
 \eps(v)K_v.
 \label{eq:correct-terminal}
\end{equation}
Conversely, replacing a terminal datum by a positive row realizing $K_v$ while
$\eps(v)=-1$ does not preserve the root marginal unless $K_v=0$ in every
observed coordinate.
\end{theorem}

\begin{proof}
Apply signed observation to the direct-sum identity at every node.  By
\eqref{eq:parity-cocycle}, the observation of the terminal source at $v$ is
$\eps(v)K_v$.  Substitution of exact terminal realizations therefore preserves
the root identity.  Since all coefficients are nonnegative and ownership is
one-use, the substituted sum is positive.

For necessity, suppose $\eps(v)=-1$ but the installed row realizes $K_v$.
Choose an observed coordinate on which $K_v\ne0$.  The correct contribution is
$-K_v$ and the installed contribution is $K_v$, changing the root coordinate
by $2K_v$.  No label bookkeeping can remove this change after the terminal
source has been replaced.
\end{proof}

\begin{corollary}[When parity-blind gluing is legal]
A terminal map proved only for canonical orientation can be applied to every
leaf only if all stopped histories are even, or if both $K_v$ and $-K_v$ have
separately proved positive realizations.
\end{corollary}

\begin{firewall}
Storing the parity in the source label is not sufficient.  The terminal
realization itself must use the label through \eqref{eq:correct-terminal}.
\end{firewall}

\section{An explicit odd-history obstruction}

\subsection{The terminal target atoms}

The canonical one-prime terminal construction uses the target kernel
\begin{equation}
 T(u)=(4\sqrt u-3)\1_{u\ge1}.
 \label{eq:T}
\end{equation}
For a terminal prime $p\ge67$, child endpoint $1\le y<67$, and small divisor
$d\mid P_{61}$, define
\begin{equation}
 t_d(p,y)=\frac1{\sqrt d}
 \left[T(py/d)-\frac1{\sqrt p}T(y/d)\right].
 \label{eq:target-atom}
\end{equation}
The canonical even and odd target masses are
\begin{equation}
 E_T(p,y)=\sum_{\substack{d\mid P_{61}\\ \mu(d)=1}}t_d(p,y),
 \qquad
 O_T(p,y)=\sum_{\substack{d\mid P_{61}\\ \mu(d)=-1}}t_d(p,y).
 \label{eq:EO}
\end{equation}
Inactive terms vanish causally.

The Target--Lorenz map used in the predecessor route is oriented from canonical
even supply to canonical odd demand.  It selects an even submeasure of target
mass $O_T$ and leaves a positive even residual.  Thus it requires $E_T\ge O_T$.

\subsection{The path \texorpdfstring{$67\to71$}{67 to 71} at endpoint 61841}

Take
\begin{equation}
 X=67\cdot71\cdot13=61841.
 \label{eq:X-obstruction}
\end{equation}
After the first rough factor $67$, the node endpoint is $923=71\cdot13$.  The
next terminal current leaf has $(p,y)=(71,13)$.  The incoming history has length
one, so Theorem~\ref{thm:parity-gluing} reverses the canonical orientation.
Actual even supply is the canonical odd mass $O_T(71,13)$, while actual odd
demand is the canonical even mass $E_T(71,13)$.

\begin{theorem}[Directed target obstruction]
\label{thm:target-obstruction}
For $(p,y)=(71,13)$,
\begin{equation}
 E_T(71,13)-O_T(71,13)>17.
 \label{eq:gap17}
\end{equation}
Consequently the reversed exact-target Hall condition
$O_T(71,13)\ge E_T(71,13)$ is impossible.
\end{theorem}

\begin{proof}
There are $239$ active squarefree divisors $d\mid P_{61}$ with $d\le923$.
The supplied verifier encloses every square root by exact dyadic rationals with
denominator $2^{192}$.  For a rational $a/b>0$, it computes
\[
 n=\left\lfloor\sqrt{\left\lfloor a2^{384}/b\right\rfloor}\right\rfloor
\]
and uses
\[
 \frac{n}{2^{192}}
 \le \sqrt{a/b}
 \le \frac{n+1}{2^{192}},
\]
with equality detected by exact integer arithmetic.  Interval addition,
multiplication, subtraction, and reciprocal are then rational operations.
Applying these operations to \eqref{eq:target-atom} and summing by M\"obius
parity yields
\begin{align*}
E_T(71,13)&\in
[238.236769968802241851178078781701\ldots,\\
&\hspace{31mm}238.236769968802241851178078781702\ldots],\\
O_T(71,13)&\in
[221.231683430583189260479623028094\ldots,\\
&\hspace{31mm}221.231683430583189260479623028095\ldots].
\end{align*}
The resulting rational lower bound for the difference exceeds $17$.
The complete rational endpoints are retained in the deterministic verification
object; the decimals above are display strings only.
\end{proof}

\begin{corollary}[Failure of parity-blind terminalization]
The canonical Target--Lorenz terminal map cannot be substituted on the leaf
\eqref{eq:X-obstruction} after history $(67)$ while preserving the complete
target datum.  Therefore the parity-blind projective composition does not prove
$c_X(j)\ge0$.
\end{corollary}

\begin{remark}
The conclusion is deliberately narrow.  It refutes the proof interface, not the
inequality $c_X(j)\ge0$.  Different rough histories may cancel globally.  Any
successful successor must retain that global cancellation or provide a new
positive joint source decomposition.
\end{remark}

\section{Fixed-depth parity resets and their obstruction}

The odd-history obstruction suggests advancing the source by an even number of
rough primes before recursing.  Source-theoretically this is exact.  Positivity,
however, imposes a second and independent constraint.

Let $p_j,p_{j+1},\ldots$ be the allowed ordered rough primes.  Let $P_j(x)$ be
the corresponding paired source, and let $b_j(x)$ be the base consisting of
histories with no allowed rough prime.  Unique least-prime factorization gives
\begin{equation}
 P_j(x)=b_j(x)\oplus
 \bigoplus_{k\ge j}p_k^{-1/2}S P_{k+1}(x/p_k).
 \label{eq:one-level}
\end{equation}

\begin{theorem}[Exact depth-$L$ expansion]
\label{thm:depth-L}
For every fixed integer $L\ge1$,
\begin{align}
 P_j(x)={}&
 \bigoplus_{r=0}^{L-1}
 \ \bigoplus_{j\le k_1<\cdots<k_r}
 (p_{k_1}\cdots p_{k_r})^{-1/2}S^r
 b_{k_r+1}\!\left(\frac{x}{p_{k_1}\cdots p_{k_r}}\right)
 \notag\\
 &\oplus
 \bigoplus_{j\le k_1<\cdots<k_L}
 (p_{k_1}\cdots p_{k_L})^{-1/2}S^L
 P_{k_L+1}\!\left(\frac{x}{p_{k_1}\cdots p_{k_L}}\right),
 \label{eq:depth-L}
\end{align}
where the $r=0$ term is $b_j(x)$.  The partition is disjoint and exhaustive.
The recursive packets have parity $(-1)^L$ and endpoint at most
$x/(p_j\cdots p_{j+L-1})$.
\end{theorem}

\begin{proof}
Classify a finite rough history by whether its length is smaller than $L$.  A
history of length $r<L$ has unique ordered primes $p_{k_1}<\cdots<p_{k_r}$ and
terminates in the corresponding base.  A history of length at least $L$ has a
unique first $L$-tuple and a remaining history in $P_{k_L+1}$.  Coefficients
factor multiplicatively and the swaps contribute $S^r$ or $S^L$.
\end{proof}

For $P=P_{61}$, write the positive finite-Euler row
\begin{equation}
 D_{P,Y}(q)=\sum_{d\mid P}\frac{\mu(d)}{\sqrt d}Q_{Y/d}(q).
 \label{eq:finite-euler-row}
\end{equation}
After signed observation, the current block in \eqref{eq:depth-L} at rough
threshold $67$ is
\begin{equation}
 B_{L,X}(q)=
 \sum_{\substack{m\le X/q\\m\ \mathrm{squarefree}\\
                   p\mid m\Rightarrow p\ge67\\\omega(m)<L}}
 \frac{\mu(m)}{\sqrt m}D_{P,X/m}(q).
 \label{eq:fixed-depth-current}
\end{equation}
For $L=2$ this is the root base minus all one-rough-prime bases.

\begin{lemma}[Finite-Euler principal term]
\label{lem:finite-euler-asymptotic}
For each fixed row $q\ge2$,
\begin{equation}
 D_{P,Y}(q)=a_q\sqrt Y+O_q(\log(2Y)),
 \qquad
 a_q=\frac{8}{q(q-1)}\prod_{\substack{\ell\le61\\\ell\ \mathrm{prime}}}
 \left(1-\frac1\ell\right)>0.
 \label{eq:finite-euler-asymptotic}
\end{equation}
\end{lemma}

\begin{proof}
The tail in \eqref{eq:Q} is a Riemann sum for
\[
 \int_1^Y t^{-1/2}\log(Y/t)\,dt=4\sqrt Y-4-2\log Y.
\]
Integral comparison for the positive decreasing summand gives
\[
 Q_Y(q)=4C_q\sqrt Y+O_q(\log(2Y)).
\]
Insert this in the finite sum \eqref{eq:finite-euler-row}.  The principal
coefficient is
\[
 4C_q\sum_{d\mid P}\frac{\mu(d)}d
 =\frac{8}{q(q-1)}\prod_{\ell\mid P}(1-\ell^{-1}).
\]
\end{proof}

We use the following standard fixed-order form of the almost-prime estimates.
It is included to make the obstruction independent of numerical evidence.

\begin{lemma}[Reciprocal rough almost-primes]
\label{lem:rough-almost-primes}
Fix $r\ge1$ and a fixed prime threshold $z$.  Put
\[
 H_z(X)=\sum_{z\le p\le X}\frac1p.
\]
For squarefree integers all of whose prime factors are at least $z$,
\begin{align}
 \sum_{\substack{m\le X\\\omega(m)=r}}
 \frac1m
 &=\frac{H_z(X)^r}{r!}+O_{r,z}(H_z(X)^{r-1}),
 \label{eq:reciprocal-almost-prime}\\
 \sum_{\substack{m\le X\\\omega(m)=r}}
 \frac{\log(2X/m)}{\sqrt m}
 &=O_{r,z}\!\left(
 \frac{\sqrt X\,H_z(X)^{r-1}}{\log X}
 \right).
 \label{eq:half-almost-prime}
\end{align}
Moreover $H_z(X)=\log\log X+O_z(1)$.
\end{lemma}

\begin{proof}
Mertens' theorem for primes gives the last assertion.  Expanding the $r$-fold
prime harmonic sum, deleting repeated primes, imposing increasing order, and
then imposing the product cutoff changes the result by
$O_{r,z}(H_z^{r-1})$; this proves \eqref{eq:reciprocal-almost-prime} by induction
on $r$.

For \eqref{eq:half-almost-prime}, the same induction gives the classical
fixed-order Landau bound
\[
 A_r(t):=\#\{m\le t:\omega(m)=r\}
 \ll_{r,z}\frac{t(\log\log t)^{r-1}}{\log t}.
\]
Partial summation against the decreasing function
$t^{-1/2}\log(2X/t)$, splitting at $\sqrt X$, yields the displayed bound.
All constants may depend on the fixed $r$ and $z$.
\end{proof}

\begin{theorem}[Fixed-depth parity-reset no-go]
\label{thm:fixed-depth-no-go}
For every fixed $L\ge2$ and fixed row $q\ge2$,
\begin{equation}
 \frac{B_{L,X}(q)}{\sqrt X}
 =a_q\frac{(-1)^{L-1}}{(L-1)!}
  (\log\log X)^{L-1}
 +O_{L,q}((\log\log X)^{L-2}).
 \label{eq:fixed-depth-asymptotic}
\end{equation}
Consequently no fixed depth can simultaneously restore canonical parity and
leave an eventually nonnegative current block:
\begin{itemize}
\item if $L$ is even, the recursive parity is canonical but
$B_{L,X}(q)<0$ for all sufficiently large $X$;
\item if $L$ is odd, the leading current term is positive but the recursive
packets retain odd parity.
\end{itemize}
\end{theorem}

\begin{proof}
Insert \eqref{eq:finite-euler-asymptotic} into
\eqref{eq:fixed-depth-current}.  The principal part is
\[
 a_q\sqrt X
 \sum_{r=0}^{L-1}(-1)^r
 \sum_{\substack{m\le X/q\\\omega(m)=r}}
 \frac1m.
\]
Lemma~\ref{lem:rough-almost-primes} shows that the term of highest degree in
$H_{67}(X)$ is the $r=L-1$ term.  The accumulated error from the
$O_q(\log(2X/m))$ remainder is, by
\eqref{eq:half-almost-prime},
$o(\sqrt X(\log\log X)^{L-1})$.  Since
$H_{67}(X)=\log\log X+O(1)$, \eqref{eq:fixed-depth-asymptotic} follows.
The parity statements come from Theorem~\ref{thm:depth-L}.
\end{proof}

The first repaired depth already fails at a moderate explicit endpoint.

\begin{theorem}[Directed two-level counterexample]
\label{thm:star-counterexample}
At $X=200000$, the depth-two current block satisfies
\begin{equation}
 B_{2,200000}(2)<-11,
 \qquad
 B_{2,200000}(3)<-2.
 \label{eq:star-counterexample}
\end{equation}
\end{theorem}

\begin{proof}
The accompanying program evaluates \eqref{eq:fixed-depth-current} with
$256$-bit MPFR intervals and directed rounding for every square root, logarithm,
addition, multiplication, division, and subtraction.  It includes all $9574$
active rough primes in row two and all $6627$ in row three.  The retained
intervals are
\begin{align*}
B_{2,200000}(2)&\in[-11.2745354467343290,-11.2745354467343289],\\
B_{2,200000}(3)&\in[-2.1151635282980810,-2.1151635282980809].
\end{align*}
The upper endpoints imply \eqref{eq:star-counterexample}.
\end{proof}

\begin{firewall}
Advancing by a fixed even number of primes repairs the sign of the recursive
packet but cannot yield a globally positive local current producer.  Increasing
the fixed depth merely postpones the same obstruction.  A viable successor must
balance parity globally, use a depth growing with the endpoint, or retain sign
inside a non-scalar positive state.
\end{firewall}

\section{A clean global parity-aware gluing theorem}

The fixed-depth no-go leaves a natural exact architecture: expand the finite
source at a fixed endpoint, retain the cumulative parity of every atom, and
solve one global common-source problem before any signed observation.

Let $E$ and $O$ be the actual positive even and odd source atoms after all
rough-history swaps, with nonnegative source coefficients $a_e$ and $b_o$.
Give atom $i$ a strictly positive target $t_i$, a score $s_i$, and a row vector
$r_i$ in a closed finite-dimensional cone $K$.  In the
conclusion-producing application one may take $K=\R_+^2$ for rows two and three,
or $K=\R_+^{65}$ for rows $2$ through $66$.

A global exact-target removal is a coefficient vector
\begin{equation}
 0\le u_e\le a_e,
 \qquad
 \sum_{e\in E}u_et_e=\sum_{o\in O}b_ot_o.
 \label{eq:global-target}
\end{equation}
Use the same $u_e$ in target, score, and every row.

\begin{theorem}[Global parity-aware common-source gluing]
\label{thm:global-gluing}
Suppose a vector $u$ satisfying \eqref{eq:global-target} also satisfies
\begin{equation}
 \sum_eu_er_e-\sum_ob_or_o\in K,
 \qquad
 \sum_eu_es_e\le\sum_ob_os_o.
 \label{eq:global-row-score}
\end{equation}
Then the signed source has the exact positive typed decomposition
\begin{equation}
 (E-O)=(\nu;B;\sigma),
 \label{eq:global-typed}
\end{equation}
where
\begin{align}
 \nu_e&=a_e-u_e\ge0,\\
 B&=\sum_eu_er_e-\sum_ob_or_o\in K,\\
 \sigma&=\sum_ob_os_o-\sum_eu_es_e\ge0.
\end{align}
Its row marginal is
\begin{equation}
 R(E)-R(O)=R(\nu)+B\in K.
 \label{eq:global-positive-row}
\end{equation}
Conversely, if the compact feasible set defined by
\eqref{eq:global-target}--\eqref{eq:global-row-score} is empty, finite-dimensional
separation supplies a dual functional that excludes every common-source
removal, not merely a chosen greedy transport.
\end{theorem}

\begin{proof}
Split each even coefficient once as $a_e=u_e+(a_e-u_e)$.  Exact target equality
uses all odd target and leaves the positive residual source $\nu$.  The two
inequalities in \eqref{eq:global-row-score} define the target-null positive row
bonus $B$ and score surplus $\sigma$, giving
\eqref{eq:global-positive-row}.  For the converse, the coefficient box is
compact and convex, the target condition is affine, and the cone constraints
are closed and convex.  If their intersection is empty, the finite-dimensional
strong separation theorem gives the asserted dual certificate.
\end{proof}

\begin{definition}[Global parity-Hall producer]
The statement $\GPHT$ asserts that the complete, parity-correct finite source at
every endpoint admits the common-source certificate of
Theorem~\ref{thm:global-gluing} with $K=\R_+^2$ for rows two and three.
\end{definition}

\begin{corollary}[Full-recursion reduction]
\label{cor:GPHT-rows}
If $\GPHT$ holds for every sufficiently large endpoint, then
$c_X(2),c_X(3)\ge0$ for every sufficiently large $X$.
\end{corollary}

\begin{proof}
The full rough-prime tree is finite at fixed $X$ and is a disjoint expansion of
the original squarefree source.  Activation and coefficient covariance are
exact by Theorem~\ref{thm:activation}.  Apply
Theorem~\ref{thm:global-gluing} once after all history parities have been
incorporated.  Equation \eqref{eq:global-positive-row} gives the two row
inequalities.
\end{proof}

\begin{remark}
The theorem is a genuine architectural repair but not a proof of $\GPHT$.
Unlike the predecessor, it does not commute a nonlinear Hall map through the
rough tree and does not reuse a root base at separate leaves.  The remaining
producer is global and finite at each endpoint; it can emit either a common
positive packing or an exact dual separator.
\end{remark}


\section{The exact two-row Mellin--Landau consumer}

This section is independent of the projective source architecture.  It proves
that eventual nonnegativity of rows two and three implies the Riemann
Hypothesis.

\subsection{Fixed-row Mellin transform}

For fixed $j\ge2$ and $\Re s>1/2$, termwise integration is absolute and
\begin{equation}
 \int_m^\infty \log\frac Xm X^{-s-1}\dd X=\frac{m^{-s}}{s^2}.
 \label{eq:ramp-mellin}
\end{equation}
With $z=s+1/2$, define
\begin{align}
 H_j(z)&=A_jj^{-z}-B_j(j+1)^{-z}
 +C_j\sum_{m\ge j+2}m^{-z},\\
 P_j(z)&=A_jj^{-z}-B_j(j+1)^{-z}
 -C_j\sum_{m=1}^{j+1}m^{-z}.
 \label{eq:Pj}
\end{align}
Then $H_j(z)=C_j\zeta(z)+P_j(z)$.

\begin{theorem}[Reciprocal-zeta row transform]
For $\Re s>1/2$,
\begin{equation}
 \boxed{
 \int_1^\infty c_X(j)X^{-s-1}\dd X
 =\frac{C_j}{s^2}
 +\frac{P_j(s+1/2)}{s^2\zeta(s+1/2)}.}
 \label{eq:mellin-transform}
\end{equation}
The right side gives meromorphic continuation to $\Re s>0$ and is analytic at
every positive real $s$.
\end{theorem}

\begin{proof}
Integrating \eqref{eq:Q} and using \eqref{eq:ramp-mellin} gives
$H_j(s+1/2)/s^2$.  In \eqref{eq:native-row}, substitute $X=kY$ and sum the
absolutely convergent Dirichlet series
\[
 \sum_{k\ge1}\frac{\mu(k)}{k^{s+1/2}}
 =\frac1{\zeta(s+1/2)}.
\]
Since $H_j=C_j\zeta+P_j$, equation \eqref{eq:mellin-transform} follows.  On the
positive real axis, $\zeta(s+1/2)$ has no zero; its pole at $s=1/2$ becomes a
zero of $1/\zeta$.
\end{proof}

\subsection{Two-row noncancellation}

For rows two and three, let $x=2^{-z}$ and $y=3^{-z}$.  Direct simplification
gives
\begin{equation}
 P_2(z)=2x-1-y,
 \qquad
 3P_3(z)=5y-x-1-3x^2.
 \label{eq:P23}
\end{equation}

\begin{theorem}[No common open-strip cancellation]
If $0<\Re z<1$, then $P_2(z)$ and $P_3(z)$ do not vanish simultaneously.
\end{theorem}

\begin{proof}
If $P_2=0$, then $y=2x-1$.  Substitution in the second equation of
\eqref{eq:P23} gives
\[
 -3x^2+9x-6=-3(x-1)(x-2)=0.
\]
But $1/2<|2^{-z}|<1$ in the open strip, excluding both $x=1$ and $x=2$.
\end{proof}

\subsection{A self-contained Landau lemma}

\begin{lemma}[Landau for a nonnegative Laplace density]
\label{lem:landau}
Let $g:[0,\infty)\to[0,\infty)$ be locally integrable and let
\[
 F(s)=\int_0^\infty g(t)e^{-st}\dd t
\]
have finite abscissa of convergence $\sigma_c$.  Then $F$ cannot be holomorphic
in a neighborhood of the real point $\sigma_c$.
\end{lemma}

\begin{proof}
Assume $F$ is holomorphic in a disk of radius $R>0$ centered at a real
$\sigma_1>\sigma_c$ with $\sigma_1-\sigma_c<R$.  Cauchy's estimate gives
\[
 \frac{|F^{(n)}(\sigma_1)|}{n!}\le MR^{-n}
\]
for some $M$.  Positivity permits differentiation under the integral and removes
absolute values:
\[
 (-1)^nF^{(n)}(\sigma_1)
 =\int_0^\infty t^ng(t)e^{-\sigma_1t}\dd t\ge0.
\]
For $0<\delta<R$, monotone convergence yields
\begin{align*}
\int_0^\infty g(t)e^{-(\sigma_1-\delta)t}\dd t
&=\sum_{n\ge0}\frac{\delta^n}{n!}
\int_0^\infty t^ng(t)e^{-\sigma_1t}\dd t\\
&\le M\sum_{n\ge0}(\delta/R)^n<\infty.
\end{align*}
Choosing $\delta>\sigma_1-\sigma_c$ contradicts the definition of
$\sigma_c$.
\end{proof}

By the change of variables $X=e^t$, the same statement applies to nonnegative
Mellin densities $f(X)X^{-s-1}\dd X$.

\begin{theorem}[Two-row positivity implies RH]
\label{thm:TRP-RH}
If $c_X(2)$ and $c_X(3)$ are nonnegative for every sufficiently large $X$, then
the Riemann Hypothesis holds.
\end{theorem}

\begin{proof}
Removing a compact initial interval changes each Mellin transform by an entire
function.  Apply Lemma~\ref{lem:landau} to the nonnegative tail.  Since the
meromorphic formula \eqref{eq:mellin-transform} is analytic at every positive
real point, the abscissa of convergence is at most zero, so the tail transform
is holomorphic in $\Re s>0$.

If $\rho$ were a nontrivial zero with $\Re\rho>1/2$, then
$s_0=\rho-1/2$ would lie in $\Re s>0$.  At least one of $P_2(\rho)$ and
$P_3(\rho)$ is nonzero, so the corresponding expression
\eqref{eq:mellin-transform} has a pole at $s_0$, a contradiction.  The
functional equation and conjugation symmetry then place every nontrivial zero
on $\Re s=1/2$.
\end{proof}

\begin{corollary}
The global parity-Hall producer $\GPHT$ implies the Riemann Hypothesis.
\end{corollary}

\section{What survives and what is withdrawn}

\begin{center}
\begin{tabular}{p{0.54\textwidth}p{0.32\textwidth}}
\toprule
Interface & Status \\
\midrule
Canonical component-row formula & exact \\
Rough-history coefficient magnitude & exact \\
Rough-history row activation & exact \\
Cumulative parity $(-1)^{|h|}$ & exact \\
Canonical Target--Lorenz terminal map & one orientation only \\
Parity-blind arbitrary-history gluing & refuted as an argument \\
Full-row nonnegativity $c_X(j)\ge0$ & open, not refuted \\
Depth-$L$ parity-aware source identity & exact \\
Fixed-depth local positive producer & impossible \\
Global producer $\GPHT$ & open and RH-bearing \\
Fixed-row reciprocal-zeta Mellin transform & exact \\
Rows two and three noncancellation & exact \\
Riemann Hypothesis & unproved \\
\bottomrule
\end{tabular}
\end{center}

The directed Target--Lorenz tail calculations from the predecessor remain valid
at their stated canonical-orientation scope.  They do not repair the sign
cocycle.  Likewise, the exact coefficient contraction and one-use source
ledgers remain useful.  The failed step is the claim that one canonical
terminal map may be substituted after arbitrary rough histories.

\section{Deterministic verification}

The accompanying verifier uses Python's arbitrary-precision integers and
\code{Fraction} class for the dyadic target obstruction, together with direct
$256$-bit MPFR intervals for the fixed-depth current block.  It checks:

\begin{enumerate}[label=(\arabic*)]
\item $16$ ordered rough histories over the fixture primes
$67,71,73,79$;
\item $309120$ exact coefficient, ratio, activation, and parity identities;
\item all $239$ active $P_{61}$ atoms in the leaf $(71,13)$;
\item the rational interval inequality \eqref{eq:gap17};
\item a disjoint and exhaustive two-level partition of every subset of the
four-prime fixture;
\item a $256$-bit directed MPFR evaluation of the complete depth-two current
block at $X=200000$;
\item the exact polynomial factorization in \eqref{eq:P23};
\item hostile mutations that drop history parity, reverse Hall capacity, or
terminalize a one-prime leaf parity-blindly.
\end{enumerate}

The retained verdict is
\begin{verbatim}
PASS_X_96500_PARITY_COVARIANT_GLUE_AUDIT
918b3ccfb7f99d764e01302ffc2f358fda0e3e2ccfc27eeb038799cf0ee05293
E_T(71,13)-O_T(71,13)>17
\end{verbatim}
The finite fixture authenticates the exact formulas and the counterexample.  It
does not prove $\GPHT$, two-row positivity, or RH.

\section{Conclusion}

The hostile reconstruction resolves the parity and activation question sharply.
Row activation and coefficient magnitude are preserved under every rough-prime
placement.  The nontrivial interface is the sign cocycle.  A single preceding
rough prime reverses the canonical one-prime leaf, and the published
Target--Lorenz map cannot be reversed even at target-mass level for the explicit
leaf $(71,13)$.

The obvious fixed-depth repair is not viable.  Odd depth leaves the recursive
parity reversed; even depth restores it but creates an eventually negative
current Euler truncation.  The exact row-two and row-three counterexample at
$X=200000$ makes this obstruction concrete.  Thus the route cannot be repaired
by merely grouping two, four, or any other fixed number of rough levels.

The clean surviving architecture is global.  At a fixed endpoint the complete
source tree is finite, so every cumulative parity can be incorporated before
one common-source transport is applied.  Theorem~\ref{thm:global-gluing} carries
all source, target, score, and row coordinates with one coefficient vector and
has an exact dual alternative.  Its two-row producer $\GPHT$ is the remaining
arithmetic theorem.  The downstream analytic chain is complete:
$\GPHT$ implies eventual positivity of rows two and three, and that positivity
implies RH through the reciprocal-zeta Mellin transforms and Landau's theorem.
No proof of $\GPHT$ is contained here.  Accordingly, RH remains unproved.

\appendix
\section{Exact interval certificate schema}

For completeness, the target certificate is determined by the following finite
algorithm.

\begin{enumerate}[label=(\roman*)]
\item Generate every squarefree divisor $d$ of $P_{61}$ with $d\le923$.
\item Determine $\mu(d)$ by parity of its prime subset.
\item Enclose each square root in \eqref{eq:target-atom} by the exact dyadic
method in Theorem~\ref{thm:target-obstruction}.
\item Apply outward rational interval arithmetic to obtain $t_d$.
\item Sum $t_d$ separately over $\mu(d)=1$ and $\mu(d)=-1$.
\item Subtract the upper odd endpoint from the lower even endpoint.
\end{enumerate}

Every operation is an integer or rational operation.  The certificate is
therefore independent of the host floating-point environment.

\section{Repository provenance and claim boundary}

The reconstruction froze the predecessor projective full-row proposal at head
\sha{20646a78c3e8843001cb49ea0c9741f6d0d446f7}, stacked on the projective
native-Lorenz proposal at
\sha{ef18bdda5a65334695008e4c1f5986833160d84f}.  The observed main branch at
the audit cutoff was
\sha{994bd4bedcc8b61cebabaf005cc69225fc3fe459}.

The present manuscript makes no claim that the full native row is negative.  It
withdraws only the parity-blind proof.  Its new theorem labels are:

\begin{center}
\begin{tabular}{ll}
\toprule
\code{L-96500} & parity-covariant projective gluing \\
\code{L-96501} & activation and coefficient covariance \\
\code{R-96500} & explicit odd-history obstruction \\
\code{L-96502} & exact depth-$L$ parity expansion \\
\code{R-96501} & fixed-depth positive-reset no-go \\
\code{L-96503} & global parity-aware common-source gluing \\
\code{T-96500} & $\GPHT\Rightarrow$ two-row positivity $\Rightarrow$ RH \\
\bottomrule
\end{tabular}
\end{center}

\begin{thebibliography}{9}
\bibitem{Riemann1859}
B. Riemann,
\emph{Ueber die Anzahl der Primzahlen unter einer gegebenen Groesse},
Monatsberichte der Berliner Akademie, 1859.

\bibitem{Landau}
E. Landau,
classical singularity theorem for Dirichlet and Laplace transforms with
nonnegative coefficients.  A self-contained form sufficient here is proved as
Lemma~\ref{lem:landau}.
\end{thebibliography}

\end{document}
